\begingroup
\renewcommand{\arraystretch}{1.45}
\newcommand{\TransscendentArticle}[1]{\par\medskip\Needspace{5\baselineskip}\pdfbookmark[1]{Art. #1}{transscendentes-ii-art-#1}\centerline{[#1.]}\nobreak\smallskip}
\fancyhead[L]{\small ZUR THEORIE DER NEUEN TRANSZENDENTEN. II.}

\Needspace{26\baselineskip}
\TransscendentArticle{3}
Zieht man auf ähnliche Weise von der Reihe
\[
\begin{gathered}\frac{1-x^{2n+2}}{1-x^{n+1}}
\\[7pt]\quad{}+\frac{x^n\,.\,1-x^{2n+4}\,.\,1-x^{n+2}}{1-x^{n+1}\,.\,1-x^{n+3}}
\\[7pt]\quad{}+\frac{x^{2n}\,.\,1-x^{2n+6}\,.\,1-x^{n+2}\,.\,1-x^{n+4}}{1-x^{n+1}\,.\,1-x^{n+3}\,.\,1-x^{n+5}}
+\text{usw.}\end{gathered}
\]
\Needspace{12\baselineskip}
die Reihe
\[
\begin{gathered}\frac{x^n\,.\,1-x^{n+2}}{1-x^{n+1}}
\\[7pt]\quad{}+\frac{x^{2n}\,.\,1-x^{n+2}\,.\,1-x^{n+4}}{1-x^{n+1}\,.\,1-x^{n+3}}
\\[7pt]\quad{}+\frac{x^{3n}\,.\,1-x^{n+2}\,.\,1-x^{n+4}\,.\,1-x^{n+6}}{1-x^{n+1}\,.\,1-x^{n+3}\,.\,1-x^{n+5}}
+\text{usw.}\end{gathered}
\]
\Needspace{16\baselineskip}
ab, so erhält man erstens
\[
\begin{gathered}\frac{1-x^n}{1-x^{n+1}}
\\[7pt]\quad{}+\frac{x^n\,.\,1-x^n\,.\,1-x^{n+2}}{1-x^{n+1}\,.\,1-x^{n+3}}
\\[7pt]\quad{}+\frac{x^{2n}\,.\,1-x^n\,.\,1-x^{n+2}\,.\,1-x^{n+4}}{1-x^{n+1}\,.\,1-x^{n+3}\,.\,1-x^{n+5}}
+\text{usw.}=\psi(x,n)\end{gathered}
\]
\Needspace{15\baselineskip}
und zweitens
\[
\begin{gathered}
1+x^{n+1}+\frac{x^{2n+3}\,.\,1-x^{n+2}}{1-x^{n+3}}
\\[7pt]\quad{}+\frac{x^{2n+3}\,.\,x^{n+2}\,.\,1-x^{n+2}\,.\,1-x^{n+4}}{1-x^{n+3}\,.\,1-x^{n+5}}+\text{usw.}\\[7pt]
=1+x^{n+1}+x^{2n+3}\psi(x,n+2)
\end{gathered}
\]
\Needspace{8\baselineskip}
Man hat also, außer im Fall \(n=0\),
\Needspace{8\baselineskip}
\[
\psi(x,n)=1+x^{n+1}+x^{2n+3}\psi(x,n+2)
\]
folglich
\[
\psi(x,n)=1+x^{n+1}+x^{2n+3}+x^{3n+6}+x^{4n+10}+\text{usw.}
\]
\Needspace{10\baselineskip}
Dagegen hat man für den Fall \(n=0\)
\[
\begin{gathered}\psi(x,0)=1+x+x^3\psi(x,2)\\[7pt]\quad{}-\frac{1-x^2\,.\,1-x^4\,.\,1-x^6}{1-x\,.\,1-x^3\,.\,1-x^5}\ \text{usw.}\end{gathered}
\]
\Needspace{8\baselineskip}
\SourceLeaf{449}
folglich
\[
\frac{1-x^2\,.\,1-x^4\,.\,1-x^6\ .\ .\ .}{1-x\,.\,1-x^3\,.\,1-x^5\ .\ .\ .}
=1+x+x^3+x^6+x^{10}+\text{usw.}
\]

\Needspace{24\baselineskip}
\TransscendentArticle{4}
Wir bezeichnen
\[
1-x\,.\,1-xx\,.\,1-x^3\ \text{usw.}\qquad\text{mit }[x]
\]
so gilt:
\[
\begin{gathered}
1+x+x^3+x^6+\text{usw.}
\\[7pt]\quad{}=\frac{1-xx\,.\,1-x^4\,.\,1-x^6\ .\ .\ .}{1-x\,.\,1-x^3\,.\,1-x^5\ .\ .\ .}\\[7pt]\quad{}=\frac{[xx]^2}{[x]}\\[7pt]
=1+x\,.\,1+xx\,.\,1+x^3\,.\,1+x^4\,.\,1+x^5\ .\ .\ .\\[7pt]\quad{}1-x^2\,.\,1-x^4\,.\,1-x^{16}\ .\ .\ .
\end{gathered}
\]
\[
\begin{gathered}1-2x+2x^4-2x^9+\text{usw.}
\\[7pt]\quad{}=\frac{1-x\,.\,1-xx\,.\,1-x^3\ .\ .\ .}{1+x\,.\,1+xx\,.\,1+x^3\ .\ .\ .}\\[7pt]\quad{}=\frac{[x]^2}{[xx]}\end{gathered}
\]
\[
\begin{gathered}
1+2x+2x^4+2x^9+\text{usw.}
\\[7pt]\quad{}=\frac{1+x\,.\,1-xx\,.\,1+x^3\ .\ .\ .}{1-x\,.\,1+xx\,.\,1-x^3\ .\ .\ .}\\[7pt]\quad{}=\frac{[xx]^5}{[x]^2[x^4]^2}\\[7pt]
=(1+x)^2(1-xx)(1+x^3)^2(1-x^4)(1+x^5)^2\ .\ .\ .
\end{gathered}
\]
\[
[-x]=\frac{[xx]^3}{[x][x^4]}
\]
\[
\begin{gathered}
(1+xy)(1+x^3y)(1+x^5y)\ .\ .\ .
\\[7pt]\quad{}\Bigl(1+\frac{x}{y}\Bigr)\Bigl(1+\frac{x^3}{y}\Bigr)\Bigl(1+\frac{x^5}{y}\Bigr)\ .\ .\ .\\[7pt]
\qquad=\frac{1}{[xx]}\Bigl\{1+x\Bigl(y+\frac{1}{y}\Bigr)+x^4\Bigl(yy+\frac{1}{yy}\Bigr)+x^9\Bigl(y^3+\frac{1}{y^3}\Bigr)+\,.\Bigr\}
\end{gathered}
\]

\Needspace{13\baselineskip}
\[
\begin{gathered}\text{usw.}+x^{(\omega-1)^2}+x^{\omega\omega}+x^{(\omega+1)^2}+\text{usw.}
\\[7pt]\quad{}=[xx]\frac{\ .\ .\ .1+x^{2\omega+3}\,.\,1+x^{2\omega+1}\,.\,1+x^{2\omega-1}\,.\,1+x^{2\omega-3}\ .\ .\ .}{x^{-\omega\omega}\,.\,x^{2\omega-1}\,.\,x^{2\omega-3}\ .\ .\ .}\end{gathered}
\]

\Needspace{23\baselineskip}
\[
[x]^3=1-3x+5x^3-7x^6+9x^{10}-\text{usw.}
\]
folgt leicht aus
\[
\begin{gathered}
\Bigl\{\Bigl(y-\frac{1}{y}\Bigr)x-\Bigl(y^3-\frac{1}{y^3}\Bigr)x^9+\Bigl(y^5-\frac{1}{y^5}\Bigr)x^{25}-\,.\,.\Bigr\}
\\[7pt]\quad{}\Bigl\{\Bigl(y+\frac{1}{y}\Bigr)x+\Bigl(y^3+\frac{1}{y^3}\Bigr)x^9+\Bigl(y^5+\frac{1}{y^5}\Bigr)x^{25}+\,.\,.\Bigr\}\\[7pt]
\qquad=\Bigl\{1-2x^8+2x^{32}-\,.\,.\Bigr\}
\\[7pt]\quad{}\Bigl\{\Bigl(yy-\frac{1}{yy}\Bigr)xx-\Bigl(y^6-\frac{1}{y^6}\Bigr)x^{18}+\text{usw.}\Bigr\}
\end{gathered}
\]
wenn man \(y=1+\omega\) setzt und daraus die Bedingungsgleichungen bildet.
\clearpage
\endgroup
